\section{Anomalies}

A quantum gauge theory can be secretly inconsistent through anomalies, subtle failures that must all cancel for the theory
to make sense at all. For the Standard Model they cancel only for one specific assignment of charges, and it is exactly
the assignment the fermions section derived.

On the substrate the anomalies cancel because the charge-and-color bookkeeping of one generation, the charges zero, one
third, two thirds, one with the color multiplicities one, three, three, one, makes every anomaly sum vanish, and the
vanishing of the hypercharge trace and the bottom-tau relation follow as the same kind of consequence. This is a strong
check, because the charges were fixed by the algebra before anyone asked about anomalies, and they then make the
anomalies cancel, which they had no obligation to do.

So the consistency of the gauge theory is forced by the same octonion structure that fixed the particle content. This is
a reproduction in one sense, the cancellation conditions are known, but it is a nontrivial one, because the charges that
satisfy them are derived rather than chosen. The remaining target is to derive the cancellation from a
discrete index theorem on the substrate from first principles, rather than confirming the known sums, which is a residual
rather than a gap.
